% RESEARCH_DRAFT_SHA256: e4266656e28a530e717fface2d9f93939296f7f7c2eb9b4c98a80bafd44e44e8
% Layout and source-reference wrappers expanded for the web renderer.
\documentclass[11pt]{article}
\usepackage[margin=25mm]{geometry}
\usepackage{fontspec,unicode-math}
\setmainfont{Latin Modern Roman}
\setmathfont{Latin Modern Math}
\usepackage{amsmath,amssymb,mathtools,xurl,hyperref}
\hypersetup{colorlinks=true,urlcolor=blue,linkcolor=blue}
\setcounter{secnumdepth}{0}
\setlength{\parindent}{0pt}
\setlength{\parskip}{5pt}
\setlength{\emergencystretch}{3em}
\newcommand{\sref}[2]{\hyperlink{source:#1:#2}{[S#2]}}
\newcommand{\cataloguescope}[1]{\paragraph{Source record.}#1}
\newcommand{\floor}[1]{\left\lfloor#1\right\rfloor}
\begin{document}
\section*{Research-attempt provenance}
This standalone research attempt retains the complete original definition body below,
including its historical source and preparation notes. Those notes describe the
earlier catalogue preparation; the new attempt and its verification are identified
separately after that body. No original ID, catalogue status, or source statement is changed.

% UnsolvedMath ID 20001202; source catalogue code AIM-COMPUTATION-0040
\section{Semialgebraic stratification of maximum likelihood}\label{20001202}
{\small UnsolvedMath ID 20001202 · Computer Science}
\subsection{Definitions and mathematical statement}

\paragraph{Shared notation.}Standard number systems are $\mathbb N,\mathbb Z,\mathbb Q,\mathbb R,\mathbb C$. Specialized notation and index ranges follow the cited source and the definitions below.


A semialgebraic set is described by finitely many polynomial equalities and inequalities combined by Boolean operations. Many finite-dimensional statistical parameter spaces are semialgebraic. Maximum-likelihood estimation optimizes a likelihood over that space, and stratification divides it into geometrically regular pieces on which rank or boundary behavior can be studied.
\paragraph{Record scope and qualifications.}
The catalogue title is retained as metadata. Named certificates, contracts, bounds or frameworks in that title are not asserted as results of the original source extract.
The question names an area of application without selecting a model, stratification, algorithm, or desired conclusion. The supplied catalogue marks this record partially solved; its present research status was not independently checked.
The source text below is rendered literally where necessary; damaged formulas and composite extracts are not silently repaired into a new mathematical claim.
\paragraph{Statement recorded in the supplied source.}
Question 26. What are the possible implications of semi-algebraic geometry to maximum likelihood estimation?

Seth Sullivant

Seth Sullivant spoke about Markov bases for hierarchical models and theoretical results about the structure of minimal Markov basis elements. 8 \href{https://aimath.org/WWN/compalgstat/compalgstat.pdf}{[S2]}
\paragraph{Source status.}
The supplied catalogue labels this source record partially solved; that metadata was checked against the supplied snapshot on 2026-10-11. The label is preserved for provenance and is not an independent verification of present open status or mathematical correctness.
\subsection{Short English statement}
The source asks about possible uses of semialgebraic geometry in likelihood estimation.
\subsection{Sources}
\begin{itemize}
\item[S1] Ulam AI and the UnsolvedMath Contributors, UnsolvedMath: A Curated Collection of Open Mathematics Problems, record 20001202 (AIM-COMPUTATION-0040). Catalogue curation and metadata: CC BY 4.0. The source extract is retained; definitions, English summaries and qualifications are editorial additions.
\url{https://huggingface.co/datasets/ulamai/UnsolvedMath}.
\item[S2] AIM Workshop Problem Lists, source linked by catalogue record AIM-COMPUTATION-0040.
\url{https://aimath.org/WWN/compalgstat/compalgstat.pdf}.
\end{itemize}
\paragraph{Preparation and review.}
The original source and UnsolvedMath contributors retain attribution for the source extract. Codex (GPT-6; exact serving revision not exposed) assisted with definitions, English summaries, scope notes and formatting. No new human mathematical authorship or independently verified solution is asserted. Permanent source ID 20001202 and code AIM-COMPUTATION-0040 checked against the complete supplied registry and current repository. Distinct upstream records, including repeated questions and broad or damaged extracts, are retained for complete coverage of the source CS category. A repeated record is not a new independent conjecture. Complete mathematical-equivalence and current-literature review has not been performed. Preparation and final review follow contribution prompts C and D, with source gaps disclosed. Structural and rendering checks do not establish a new theorem or open status.
\label{end:20001202}

\section{Literature and context check}
The original AIM Computational Algebraic Statistics
 workshop PDF was consulted. The preserved extract does not specify every model,
 parameterization or performance criterion needed by its broad question. The
 following declared model is an actual test of one route, not a replacement
 target or a claim of novelty for elementary statistical or algebraic facts.
\section{Research attempt: A semialgebraic likelihood graph still needs boundary and fiber analysis}
Consider Bernoulli data with $n_1$ successes and $n_0$ failures, nonnegative
integers with $N=n_0+n_1>0$. On the closed interval $0\le p\le1$ the likelihood
is $p^{n_1}(1-p)^{n_0}$, interpreting a zero exponent as one. If both counts
are positive, its log derivative is $(n_1-Np)/(p(1-p))$, changing sign exactly
at $p=n_1/N$. If one count is zero, the likelihood is monotone and its unique
maximizer is the corresponding endpoint. Thus the closed-model MLE graph is
the semialgebraic relation $Np=n_1$, $0\le p\le1$ with the count constraints.
For the open interval, either zero-count case instead has no maximizer: its
supremum is approached only at an excluded endpoint. An algebraic score equation
alone does not record that distinction.

Now parameterize the same closed family by $p=\theta^2$, $-1\le\theta\le1$.
For $0<n_1<N$ its MLE set is the two points
$\theta=\pm\sqrt{n_1/N}$. They yield one fitted distribution. For $n_1=0$
the set is $\{0\}$, and for $n_0=0$ it is $\{-1,1\}$. This proves an exact
parameter-fiber stratification in a small model, including endpoint cases.
It shows why a count of parameter maximizers cannot by itself count distinct
statistical explanations.

As a further audit, take $n_0=n_1=1$. The polynomial likelihood is
$\theta^2(1-\theta^2)$, with derivative $2\theta(1-2\theta^2)$.
Its stationary point zero is a strict local minimum, whereas
$\pm1/\sqrt2$ are the maxima. Listing all roots of the polynomial score
therefore does not solve the likelihood optimization even in degree four.
The second derivative at zero is two, verifying the rejected candidate.

The concrete route works here because monotonicity between finitely many
critical points is explicit. No general efficient stratification algorithm or
complexity guarantee follows for arbitrary semialgebraic statistical models.
\section{Verification and adversarial audit}
All zero-count cases and the excluded-boundary version were checked.
The folded parameterization was kept distinct from its observable distribution,
and the nonmaximizing stationary point was explicitly classified.
\paragraph{Audit corrections.} Score roots, parameter fibers, fitted distributions and attainment were separated; no graph-independent general MLE theorem was claimed.
\paragraph{External computation.} None; symbolic arguments and boundary checks only.
\section{Scope and remaining gaps}
The original asks for possible implications across models. The first gap
for an algorithmic extension is globally comparing strata and boundary limits
with a justified complexity bound, rather than only solving stationary equations.
\paragraph{Integration review.} Codex (GPT-6; exact serving revision not exposed) reviewed the written arguments and expanded layout and source-reference wrappers. Repository metadata and the production web renderer were checked using Python 3.13 and JavaScript with MathJax 3.2.2. This model-assisted review is not independent human proof verification; no PDF compilation or formal proof checking is asserted.

\section{Completion Estimate}
\noindent\textbf{COMPLETION: 0\%}
The calculations below the preserved source statement investigate a concrete route,
but establish no substantiated new progress toward the entire original target.
Known special cases and failed sufficient conditions are not credited as a solution.
\section{References}
The original source references and catalogue attribution are preserved in the
Sources section above. The following URLs were actually consulted for this attempt;
retention of an additional original reference does not claim fresh consultation of it.
\begin{itemize}
\item \url{https://aimath.org/WWN/compalgstat/compalgstat.pdf} (consulted 2026-10-11). 
\end{itemize}
\section{Model and contributor disclosure}
OpenAI gpt-6-astra, requested reasoning configuration ultra, generated this research
attempt on 2026-10-11. The exact serving revision is not exposed. Codex GPT-6
performed orchestration, repository-format review, and a model-assisted mathematical
audit of the written arguments; its exact serving revision
is also not exposed. The mathematical route, boundary checks and final adversarial
self-audit followed contribution prompts A then B in
\texttt{docs/contribution\_prompts.md}, after reading \texttt{CONTRIBUTING.md}.
Python 3.14.2 assembled the standalone documents; any mathematical computation is
specified in Verification. No theorem prover was used. No human mathematical
authorship, substantial human proof correction, or independent human proof checking
is asserted. Model-assisted review, source retrieval, and structural checks do not
establish correctness or novelty. Full private reasoning transcripts are unavailable
as reproducibility artifacts; the complete arguments, inputs and outcomes actually
used are stated here or in the companion batch audit.
\section{Final status}
Unresolved as a new research contribution. This is a concrete route investigation,
not a claim of new complete resolution or a current-status certification. A known
answer reproduced for context does not become a new contribution. The catalogue's
source-status label is preserved; this attempt status does not assert that a
source-labelled solved question is open.
\end{document}
